\properlane{purification-and-rest}{Christ gives rest through humility, purity, and peace}

\subsection{Refreshment and the cross}
Christ invites all who labor and are burdened to come to him for
refreshment. In the Gospel he blesses mourners, people hungering for
justice, and those persecuted for it. The Offertory adds righteous people
whose bodies have been killed. Refreshment understood as exemption from
these afflictions would contradict the other texts heard in the same Mass.
Chrysostom confronts that difficulty within Matthew 11 itself, where the
invitation is followed by the easy yoke and the light burden.

The following verses explain what disciples are to learn: Christ's
meekness and humility of heart. Chrysostom explicitly connects that
humility with the beginning of the Beatitudes. He also asks how the light
burden agrees with the narrow gate and the cross. His answer includes
Christ's example, the reward promised, and a comparison between the
burden of obedience and the burdens of sin. These contextual verses
explain the single appointed Alleluia verse; they are not additions to
its words.\footnote{Matt.~11:25--30; Chrysostom,
\work{Matthew}, XXXVIII.3--5.}

The comparison concerns experience as well as future reward. Chrysostom
describes the anxious effort to preserve possessions and the anger that
keeps a quarrel alive. Refusal to retaliate can end a contest which
retaliation would prolong. Humility removes the pride which continually
demands satisfaction. He thus gives a reason why obedience can afford
present relief, even when it requires restraint or loss. The rest he
describes is related to the changed will of the disciple.

He also invokes the suffering of apostles and martyrs and the greater
glory to come. His explanation retains actual toil. Both the yoke and its
ease are spoken of, he says, so that disciples will neither flee it as
unbearable nor despise it as requiring no effort. A person can therefore
be refreshed by Christ while still struggling to act justly. The
Sunday Collect's confessed entanglement in sins can be compared with
the bondage Chrysostom says Christ removes. Its request for release
also makes moral formation a petition for divine help.

\subsection{How humility changes what a person desires}
Augustine begins his exposition of the Beatitudes with poverty of spirit
as humility and fear of God. Pride is swollen self-importance; the humble
person submits to divine authority and can learn. This is why humility
belongs at the beginning of his account. The person who refuses instruction
cannot be corrected by the later teaching. The Introit's upright praise,
understood by Augustine as submission to God, describes a compatible
disposition at the opening of the Mass.\footnote{Augustine,
\work{Sermon on the Mount}, I.1.3; Ps.~32, enarratio I, \S1.}

In Augustine's sequence, meekness includes receptiveness to Scripture,
especially when a passage is not yet understood. Mourning then follows
the recognition of attachments which held the soul in bondage. His
account concerns the sorrow of conversion: former delights have been
lost, while love of eternal goods is still being formed. Hunger and
thirst for justice name the effort to desire the good strongly enough
to leave the old attachments. The pain of relinquishing a familiar
pleasure explains why the sequence contains both sorrow and hope.
\footnote{Augustine, \work{Sermon on the Mount}, I.2.5--6; I.3.10.}

Mercy toward others is the next act in this formation. A person seeking
help from one stronger should help another who is weaker. Augustine's
reasoning puts dependence and generosity together: the need to receive
assistance is compatible with the duty to give it. The feast's Collect
asks mercy amid many intercessors. Although Augustine does not expound
that prayer, his account of mercy explains a fitting disposition for
those who ask it: willingness to relieve another's misery.
\footnote{Augustine, \work{Sermon on the Mount}, I.2.7; I.3.10.}

In the Lesson, the servants of God receive the seal before John sees
the multitude praising God and the Lamb. Those standing before the
throne wear white robes and hold palms. Bede understands the robes
as baptism and the palms as the triumph of the Cross: they have
overcome the world in Christ. At the praise which concludes the
appointed passage, he explains that the Church acknowledges in
each member what it has received from God.\footnote{Apoc.~7:2--3,9--12;
Bede, \work{Explanation}, I, on Apoc.~7:9,12, pp.~52--53.}

Read beside Christ's invitation to the burdened, the baptismal
robes suggest that the labor of conversion begins with a gift
from God. Relinquishing sinful attachments is a response to
holiness received. The disciple needs Christ's help both to
begin and to persevere in the meekness and mercy Augustine describes.

The Sunday Secret asks for growth in service and for God to complete what
he granted without merit on the recipients' part. Schuster interprets
the requested service through ministry at the altar. An increase of
service can accordingly be understood as part of the formation being
asked for. The feast's Secret likewise asks that the offerings be
saving for those who make them, through God's mercy. Offering and
continued need remain present in both prayers.\footnote{Schuster,
\work{Sacramentary}, III, p.~188; MR 1962, pp.~417, 703, Secrets.}

\subsection{The good which hunger does not take away}
The Gradual's promise of no want is harder to understand when heard
beside the blessing of hunger. Augustine faces the devout poor person's
objection directly. A rich wrongdoer may die surrounded by costly goods,
while the person attending church and seeking God remains without
them. The difference is observable; an explanation must say what
good the promise concerns.

Augustine describes a heart rich in justice, truth, charity, faith,
and endurance. These are goods the rich sinner can lack while possessing
an ivory bed. He then names Christ as living bread and cites the
Beatitude promising fullness to those who hunger and thirst for justice.
The apparent contradiction is answered by the object of desire.
The hungering disciple seeks a good whose possession differs from
being materially prosperous. This is an explicit connection between
the Gradual's Psalm and the Gospel in Augustine's exposition.
\footnote{Augustine, Ps.~33, enarratio II, \S\S14--15;
NPNF I.8, Psalm XXXIV, \S\S13--14.}

Aquinas gives several accounts rather than reducing every deprivation
to one explanation. A person can possess spiritual treasure, be
content with enough, receive providential assistance, or be treated
like a patient from whom a physician withholds a desired food.
The last comparison explains why a request's fulfillment can depend
on what benefits the recipient. These alternatives prevent the
Gradual from being heard as a promise that a holy person will always
obtain the material thing desired. They also leave room to ask God
for actual help in bodily need.\footnote{Aquinas,
\work{Super Psalmos}, Ps.~33, nn.~10--11.}

Chrysostom offers a different emphasis in his explanation of the meek
inheriting the earth. Augustine understands an enduring inheritance,
the land of the living where the saints rest. Chrysostom argues for
a sensible present reward: meekness can preserve possessions which
anger and rashness would destroy. He also retains future gifts.
\footnote{Augustine, \work{Sermon on the Mount},
I.2.4; Chrysostom, \work{Matthew}, XV.5.}

\subsection{A pure heart and active peacemaking}
For Augustine, the pure heart is single. The comparison with sight is
precise: bodily light requires healthy eyes, and the sight of God
requires the inward faculty appropriate to him to be purified.
Singleness concerns the will and attention directed toward God.
His explanation follows the earlier relinquishment of contrary
attachments; it gives a reason why desire and vision belong within
the same account of conversion.\footnote{Augustine,
\work{Sermon on the Mount}, I.2.8.}

Chrysostom's exposition specifies a further practical difficulty.
Generous almsgiving does not make impurity irrelevant. His example
is the person who gives freely and avoids greed but still sins
against chastity. The blessing of purity asks for a life whose
different actions are answerable to God. Read beside Augustine's
singleness, this excludes treating one conspicuous virtue as
permission to retain an incompatible habit.\footnote{Chrysostom,
\work{Matthew}, XV.6.}

Augustine next explains peace as an order within the person.
Desires are subject to reason, and reason is subject to the truth,
the only-begotten Son of God. The highest human faculty itself
needs submission; its rule of what is lower depends on obedience
to what is higher. This explains how the peacemaker resembles
God and why peace can survive external opposition. A hostile
act need not make the disciple's desires rebel against the Son.
\footnote{Augustine, \work{Sermon on the Mount}, I.2.9.}

Chrysostom gives peacemaking an outward task: bringing others who
are divided into reconciliation. He connects the title of God's
children with the work of the only-begotten Son in uniting the
alienated. He then reads the persecution blessing as a necessary
qualification. An agreement which requires wrongdoing is not
blessed peace. The two expositions thus give distinct, compatible
accounts of peace: order under Christ within the person, and
reconciliation in the person's dealings with others.
\footnote{Chrysostom, \work{Matthew}, XV.7.}

The Communion repeats purity, peacemaking, and persecution for
justice in that order. Its short selection returns the hearers
to the promised sight of God, likeness to the Father, and kingdom
under opposition. In this interpretation, reception of the
sacrament is accompanied by the continued conversion these
promises describe. The Postcommunion requests for protection
are fitting because moral formation remains exposed to danger.

\subsection{Persecution, peace, and the body's resurrection}
The Gospel qualifies persecution: it is for justice. Reviling is
qualified further by falsehood and its being for Christ.
Chrysostom expressly rejects the supposition that being accused
by itself makes a person blessed. This keeps fidelity from being
confused with any quarrel in which a Christian happens to be
involved. Reconciliation remains a duty, and agreement in evil
remains excluded.\footnote{Chrysostom, \work{Matthew}, XV.7.}

Augustine explains the return of the kingdom promise in the eighth
Beatitude as the showing of a completed formation. Persecution
tests what has been established within. His claim concerns the
strength of the person's adherence to God, not the harmlessness
of the persecutor's acts. The Offertory explicitly acknowledges
bodily death while proclaiming peace in God's hand. Its words
permit a distinction between what violence does to the body
and whether the righteous person is abandoned by God.
\footnote{Augustine, \work{Sermon on the Mount}, I.2.9--I.3.10;
Wis.~3:1--3.}

Gregory develops an additional consequence of the righteous person's
suffering while explaining Job's image of a tree which can grow
again after being cut down. In \work{Moralia} XII.5.6 he directly
quotes Wisdom 3:2: the unwise judge the righteous person's death
to be misery. He interprets the despised preaching and killed
body as the apparently dried root and trunk. Through the Spirit,
the example awakens virtue in other people. His account concerns
fruit in the living, arising from the sufferer's fidelity.
\footnote{Gregory the Great, \work{Moralia in Iob}, XII.5.6,
on Job~14:8--9, directly invoking Wis.~3:2.}

Gregory's following paragraph gives a distinct interpretation of
the tree's renewed growth as bodily resurrection. These are two
explanations, with different subjects of renewal: those influenced
by the just person's example, and the just person's own body.
Both preserve what the unwise fail to perceive when they suppose
death has ended the matter. His commentary is on Job and reuses
Wisdom; its application to the Offertory follows that actual
reuse.\footnote{Gregory, \work{Moralia}, XII.5.7.}

The Common Preface describes the life toward which Christ's
refreshment leads as thanksgiving with the heavenly company.
The saints' rest includes praise. The feast's Postcommunion asks
that veneration continue in joy and that the saints continually
pray for the faithful. The Sunday's petition then names the
continuing danger of those who have shared in divine gifts.
Schuster explains protection as grace to stand within trial.
This makes the closing request compatible with the Gospel's
blessed perseverance and Chrysostom's rest under the cross.
\footnote{Schuster, \work{Sacramentary}, III, p.~189;
V, p.~211.}

\subsection{Four senses}
\begin{fourSenses}
\item[Literal.] Jesus invites the burdened and blesses humility,
mercy, purity, peace, and fidelity under righteous persecution.
The prayers ask pardon, fruitful offerings, continued service,
and protection; the Offertory proclaims the peace of the just.
\item[Allegorical.] The meek Son gives rest and reconciles those
divided. Through him, God's children give thanks with the
heavenly worshippers. His victory permits suffering and death
to be understood in relation to life with him.
\item[Moral.] Humility makes a person teachable; desire is
redirected toward justice; mercy assists the weak; purity
requires singleness and chastity; peace requires ordered
desires and practical reconciliation. Opposition tests
fidelity under the Gospel's own qualifications.
\item[Anagogical.] Present goods anticipate the promised sight
of God and the kingdom. The righteous dead's peace and the
resurrection still awaited distinguish present refreshment
from the final transformation which Augustine says exceeds
words.
\end{fourSenses}
